% !PRINTED 455
\weberpage{455}

%<<CONT
and the invariability under $\chi$, that besides this one term the two corresponding ones

\[
x_2^{h_1} x_3^{h_2} x_1^{h_3}, \quad x_3^{h_1} x_1^{h_2} x_2^{h_3}
\]

also occur in the function, unless $h_1 = h_2 = h_3$. To this is added the condition of invariability under $\omega$. We now wish to see how we can satisfy these requirements, and first of all in such a way that the degree $m$ of the invariant, that is, the sum $h_1 + h_2 + h_3$, becomes as small as possible.

The condition (1) evidently cannot be satisfied if $m < 3$. If this sum is $= 3$, then $h_1 = h_2 = h_3 = 1$ must hold; but the product $x_1 x_2 x_3$ is evidently not unchanged under $\omega$. The smallest value that comes into consideration is therefore $m = 4$, and we therefore have to seek out all non-negative solutions of

\[
\begin{aligned}
h_1 + h_2 + h_3 &= 4 \\
h_1 + 2h_2 + 4h_3 &\equiv 0\ (\mathrm{mod}\ 7)
\end{aligned}
\]
%CONT>>

If we eliminate $h_1$, it follows that

\[
h_2 + 3h_3 \equiv 3\ (\mathrm{mod}\ 7),
\]

and from this result the only possible solutions:

\[
\begin{aligned}
h_3 &= 0, & h_2 &= 3, & h_1 &= 1 \\
h_3 &= 1, & h_2 &= 0, & h_1 &= 3 \\
h_3 &= 3, & h_2 &= 1, & h_1 &= 0,
\end{aligned}
\]

and consequently the \gesperrt{only} form of degree $4$ unchanged by $\chi$ and $\tau$

\weq{2}{f\,(x_1, x_2, x_3) = x_1^3 x_3 + x_2^3 x_1 + x_3^3 x_2.}

In order to test the influence of the substitution $\omega$ on the function $f$, we set

\weq{3}{
\begin{aligned}
x_1 &= \alpha y_1 + \beta y_2 + \gamma y_3 \\
x_2 &= \beta y_1 + \gamma y_2 + \alpha y_3 \\
x_3 &= \gamma y_1 + \alpha y_2 + \beta y_3,
\end{aligned}
}

in which the coefficients $\alpha, \beta, \gamma$ have the meaning fixed in \S.\,115, and assume that by (3) the transformation

\weq{4}{F\,(y_1, y_2, y_3) = f\,(x_1, x_2, x_3)}

is accomplished. We form the second derivatives of $F$ with respect to $y_1, y_2, y_3$, with regard to (3) and to the formulae

\[
\begin{aligned}
&\tfrac{1}{6} f''(x_1, x_1) = x_1 x_3, \quad \tfrac{1}{6} f''(x_2, x_2) = x_2 x_1, \quad \tfrac{1}{6} f''(x_3, x_3) = x_3 x_2 \\
&\tfrac{1}{3} f''(x_2, x_3) = x_3^2, \quad \tfrac{1}{3} f''(x_3, x_1) = x_1^2, \quad \tfrac{1}{3} f''(x_1, x_2) = x_2^2.
\end{aligned}
\]
